\documentclass{article}
\usepackage{graphicx} % Required for inserting images

\title{Implementation of dielectric ELC in ESPResSo \\ 
\large Bachelor's Thesis}

\author{Julien Hemminger \\ 
\small Simulation Technology B.Sc.}
\date{April 2026}


\begin{document}

\maketitle

\section{Motivation}
why a new elcic impl?
there already is a elc impl, which can handle diel interfaces in espresso but:
* has error with ELCIC, nobody understands why
* unmaintainable code, hard to read
so the goal is to replace the old elc impl in espresso with a new one

\section{Prototyping in Python}
in python: things are easier/faster to develop, test, prototype, visualize, validate,...

\subsection{Regular ELC}
i use iterative/test-driven developemnt (and break the problem down into smaller problems)
* systems i look at
    * big box(negligable pbc influence) / if we make the box sizes $l_x$, $l_y$ very large, the influence of periodic images in x and y direction diminishes
    * small, neutral box(do non neutral correction yet)
    * full range of parameters (madelung, random params, extreme params, etc.)


* for each system i do
    * compute the -usually slow, expensive, limited- analytical/brute-force/reference solution(energy, force)
    * adjust my custom elc to fit the analytical solution
    * do parameter sweep
\begin{figure}
    \centering
    \includegraphics[width=1\linewidth]{mar3_regular_elc_done.png}
    \caption{Enter Caption}
    \label{fig:placeholder}
\end{figure}


\subsection{Dielectric ELC}
\subsubsection{Single Plate}
* use "ewald sum" as reference
* works for most params, but there are exceptions
\begin{figure}
    \centering
    \includegraphics[width=1\linewidth]{apr27_single_plate_elcic_mostly_works_but_some_outliers.png}
    \caption{Enter Caption}
    \label{fig:placeholder}
\end{figure}

* TODO talk about future plans, problems, etc.

* CURRENTLY HERE

\subsubsection{Dual Plates}
* couldnt get anasol to converge
\begin{figure}
    \centering
    \includegraphics[width=1\linewidth]{apr9_cant_get_dual_plate_ana_to_converge.png}
    \caption{Enter Caption}
    \label{fig:placeholder}
\end{figure}


\section{Implementing in ESPResSo}
* c++ impl: upcoming
with kokkos support


\section{NOTES TO SELF}
* TODO
    * how detailed do i really go?
    * formulas - YES but for what?
    * code snippets? NO, hard to look at, not fun to read
    * what other plots do i need?
* python plots
    * pay attention to labels, etc
        * e.g. my "analytical" solutions are sometimes actually analytical (U=1/r'2) and other times numerical (ewald oder so)


\end{document}
